\begin{topic}[id=service_robots,
             difficulty={Mittel},
             type={Anwendungsanalyse + Bewertung},
             prerequisites={Grundlagen Robotik; Interesse an Interaktion und Einsatzszenarien},
             practice={optional},
             owner={Dozententeam},
             updated={2026-04-09},
             tags={ Service-Roboter, HRI, Sicherheit, Navigation, Ethik, Normen }]{Service-Roboter}

\TopicDescription{Service-Roboter operieren in menschennahen Umgebungen: Navigation, Interaktion (HRI), Sicherheit, Datenschutz und Akzeptanz sind Schlüsselfaktoren. Wir diskutieren typische Architekturbausteine, Sensorik, Normen/Leitlinien und Praxisgrenzen (Betriebssicherheit, Wartung). Ziel ist eine nüchterne, ganzheitliche Einordnung anwendungsnaher Servicerobotik.}

\TopicLeitfragen{\item Welche HRI-Designprinzipien sind praxistauglich?
\item Wie balanciert man Sicherheit, Datenschutz und Nutzbarkeit?
\item Welche Einsatzfelder sind heute wirklich reif?
}
\TopicScope{\item Keine Produktvergleiche.
\item Keine Rechtsberatung.
}
\TopicPractice{Optionale Praxisidee: Service-Szenario skizzieren (HRI-Fluss + Sicherheitskonzept).}

\TopicKeywords{Service-Roboter, HRI, Sicherheit, Navigation, Ethik, Normen}

\nocite{iso13482Overview,hriSurvey}
\end{topic}
